\chapter{Dari Persamaan ke Solver}
\label{ch:numerik}

Setelah persamaan-persamaan atur, kuliah metode numerik fluida bergerak ke cara menyelesaikannya:
metode volume hingga, kopling tekanan dan kecepatan, turbulensi dalam tiga pertemuan, perpindahan
panas dan spesies, aliran multifase, dan satu pertemuan khusus tentang recurrence CFD. Bab ini
merangkum semuanya kecuali yang terakhir, yang mendapat babnya sendiri di \cref{ch:rcfd}.

\section{Metode volume hingga}

Metode volume hingga membagi domain menjadi sel-sel kecil dan menegakkan neraca \cref{eq:transport}
untuk setiap sel. Mengintegrasikan persamaan atas volume sebuah sel, lalu memakai teorema divergensi
Gauss, mengubah integral divergensi menjadi jumlah fluks yang melewati sisi-sisi sel:
\begin{equation}
\frac{\dd}{\dd t}\int_V\rho\phi\dd V+\sum_{\text{sisi }f}\big(\rho\vu\phi-\Gamma\nabla\phi\big)_f\cdot\bm n_fA_f=\int_VS_\phi\dd V.
\end{equation}
Keindahan bentuk ini adalah sifat konservatifnya. Fluks yang keluar dari satu sel lewat sebuah sisi
adalah persis fluks yang masuk ke sel tetangganya lewat sisi yang sama, sehingga tidak ada massa,
momentum, atau energi yang hilang atau muncul di dalam domain kecuali lewat batas dan sumber. Di
\cref{ch:gdl} kita melihat bahwa bentuk ini adalah message passing: setiap sel mengumpulkan ``pesan''
fluks dari tetangganya.

Nilai di sisi sel harus diperkirakan dari nilai di pusat sel. Untuk suku difusi, beda pusat bekerja
baik. Untuk suku konveksi, pilihannya lebih rumit. Beda pusat akurat orde dua, tetapi menghasilkan
osilasi yang tidak fisis kalau konveksi jauh lebih kuat daripada difusi di skala sel, yaitu kalau
bilangan Péclet sel lebih besar dari dua. Skema \istilah{upwind} mengambil nilai dari sel di hulu
aliran. Ia stabil dan tidak berosilasi, tetapi hanya orde satu, dan galatnya berperilaku seperti
difusi tambahan yang meratakan gradien tajam, gejala yang disebut \istilah{numerical diffusion}. Skema
orde lebih tinggi seperti second-order upwind dan skema dengan limiter mencari jalan tengah. Kuliah
memperlihatkan contohnya pada perambatan gelombang permukaan, di mana antarmuka air dan udara
mengabur dengan skema orde satu dan tetap tajam dengan skema yang lebih baik.

Untuk waktu, metode eksplisit menghitung keadaan baru langsung dari keadaan lama. Murah per langkah,
tetapi langkahnya dibatasi oleh syarat Courant--Friedrichs--Lewy \citep{courant1928},
\begin{equation}
\mathrm{CFL}=\frac{u\,\Delta t}{\Delta x}\le C_{\max},
\end{equation}
yang secara intuitif berarti informasi tidak boleh melompati lebih dari sekitar satu sel dalam satu
langkah. Metode implisit menyelesaikan sistem persamaan untuk keadaan baru. Mahal per langkah,
tetapi stabil untuk langkah yang jauh lebih besar.

\section{Tekanan dan kecepatan}

Untuk aliran tak termampatkan, tekanan tidak punya persamaannya sendiri. Ia harus diatur sedemikian
rupa sehingga medan kecepatan yang dihasilkan persamaan momentum memenuhi kekekalan massa.
Algoritme SIMPLE \citep{patankar1972,patankar1980}, yang dibahas langkah demi langkah di kuliah,
menyelesaikan masalah ini secara iteratif. Pertama, tebak medan tekanan dan selesaikan persamaan
momentum untuk mendapat kecepatan sementara. Kecepatan ini umumnya belum memenuhi kekekalan massa.
Kedua, tulis hubungan antara koreksi kecepatan dan koreksi tekanan. Ketiga, masukkan hubungan itu
ke persamaan kekekalan massa, sehingga diperoleh persamaan untuk koreksi tekanan. Selesaikan,
perbarui tekanan dan kecepatan, lalu ulangi sampai konvergen. Pendekatan yang menyelesaikan
persamaan-persamaan ini bergantian disebut \istilah{segregated solver}, sedangkan solver
\istilah{coupled} menyelesaikannya bersama dalam satu sistem besar.

Kuliah menekankan perbedaan antara konvergensi numerik dan konvergensi fisis. Residu yang turun
beberapa orde menunjukkan bahwa sistem persamaan diskret sudah terselesaikan. Itu belum berarti
besaran yang kita pedulikan, misalnya gaya hambat atau temperatur di sebuah titik, sudah tidak
berubah lagi, apalagi bahwa hasilnya benar secara fisis. Memantau besaran fisis dan menguji
kebebasan hasil terhadap ukuran grid adalah bagian wajib setiap simulasi.

\section{Turbulensi}

\subsection{Rata-rata Reynolds}

Aliran turbulen penuh dengan fluktuasi acak di semua skala. Untuk teknik, biasanya cukup mengetahui
rata-ratanya. Dekomposisi Reynolds menulis setiap besaran sebagai rata-rata ditambah fluktuasi,
$\vu=\bar{\vu}+\vu'$. Ketika dekomposisi ini dimasukkan ke persamaan Navier--Stokes dan dirata-ratakan,
hampir semuanya rapi kecuali satu suku baru dari nonlinearitas, yaitu \istilah{Reynolds stress}
$-\rho\,\overline{u_i'u_j'}$. Suku ini tidak diketahui dari besaran rata-rata, sehingga persamaannya
tidak tertutup. Inilah \istilah{closure problem} turbulensi.

Boussinesq mengusulkan agar Reynolds stress diperlakukan seperti tegangan kental, dengan sebuah
viskositas turbulen $\mu_t$ yang jauh lebih besar dari viskositas molekuler. Masalahnya bergeser
menjadi menentukan $\mu_t$. Model $k$-$\varepsilon$ \citep{launder1974} menyelesaikan dua persamaan
transport tambahan, untuk energi kinetik turbulen $k$ dan laju disipasinya $\varepsilon$, lalu menulis
$\mu_t=\rho C_\mu k^2/\varepsilon$. Konstanta-konstantanya ditetapkan dari aliran sederhana yang
dipahami dengan baik, seperti peluruhan turbulensi di belakang grid dan penyebaran jet. Model
$k$-$\omega$ SST \citep{menter1994} memakai formulasi $k$-$\omega$ di dekat dinding dan beralih ke
$k$-$\varepsilon$ di aliran bebas, dan menjadi salah satu model paling banyak dipakai di industri.
Ketika anisotropi turbulensi penting, seperti di dalam siklon gas atau pusaran ujung sayap, model
Reynolds stress menyelesaikan persamaan transport untuk setiap komponen tegangan.

Turbulensi sering digambarkan sebagai kaskade energi. Pusaran besar menerima energi dari aliran
rata-rata, pecah menjadi pusaran yang lebih kecil, dan seterusnya, sampai pusaran terkecil cukup
kecil untuk diredam viskositas menjadi panas. Seperti air terjun bertingkat, energi mengalir dari
kolam besar ke kolam-kolam yang semakin kecil, dan model turbulensi hanya perlu tahu berapa banyak
air yang jatuh ke dasar, bukan bentuk setiap riaknya.

\subsection{Large eddy simulation}

\istilah{Large eddy simulation} (LES) menyelesaikan pusaran besar secara langsung dan hanya memodelkan
pusaran yang lebih kecil dari ukuran grid. Model Smagorinsky \citep{smagorinsky1963} memberi viskositas
sub-grid $\mu_{\mathrm{SGS}}=\rho(C_s\Delta)^2|\bar{\bm S}|$, sebanding dengan kuadrat ukuran grid dan
laju deformasi yang terselesaikan. LES jauh lebih mahal dari RANS karena harus berjalan dalam waktu
dengan grid halus, tetapi jauh lebih murah dari DNS, dan ia menangkap struktur aliran tak tunak
yang hilang dalam rata-rata.

\section{Panas, spesies, dan reaksi}

Dalam aliran turbulen, panas dan zat terlarut juga diaduk oleh pusaran. Konduktivitas efektif
menjadi jumlah konduktivitas molekuler dan sumbangan turbulen,
$\lambda_{\mathrm{eff}}=\lambda+\mu_tc_p/\mathrm{Pr}_t$, dengan bilangan Prandtl turbulen biasanya sekitar
0,85. Radiasi menambah suku sumber dalam persamaan energi, dan reaksi kimia menambah suku sumber
dalam persamaan spesies dan energi sekaligus. Kuliah menyinggung contoh-contoh industri, dari nosel
Laval sampai pemanasan logam cair, untuk menunjukkan bahwa sebagian besar kerja pemodelan ada pada
suku-suku sumber ini, bukan pada solvernya.

\section{Aliran multifase}

Banyak aliran industri melibatkan lebih dari satu fase: gelembung gas di dalam cairan, butiran di
dalam gas, tetesan di dalam udara. Kuliah membahas beberapa keluarga pendekatan.

Metode \istilah{volume of fluid} \citep{hirt1981} melacak antarmuka antara dua fluida yang tidak
bercampur dengan sebuah fraksi volume yang bernilai satu di satu fase dan nol di fase lain. Massa jenis
dan viskositas dirata-ratakan menurut fraksi itu, dan tegangan permukaan masuk lewat lompatan
tekanan Young--Laplace di antarmuka, $\Delta p=\sigma\kappa$, dengan $\kappa$ kelengkungan antarmuka.
Contohnya di kuliah termasuk gelombang di tangki yang bergoyang dan aliran kapiler di saluran mikro.

Untuk partikel atau tetesan yang tersebar, pendekatan Euler--Lagrange melacak setiap partikel, atau
sekelompok partikel yang mewakilinya, dalam kerangka Lagrange, sementara fase kontinu diselesaikan di
grid. Kopling bisa satu arah, ketika partikel hanya dibawa aliran, atau dua arah, ketika partikel
juga memberi gaya balik ke fluida. Untuk aliran butiran padat, metode elemen diskret \citep{cundall1979}
menyelesaikan gerak setiap butir dan tumbukannya, dan gabungannya dengan CFD dipakai untuk unggun
terfluidisasi dan pengangkutan pneumatik. Ketika partikel terlalu banyak untuk dilacak satu per satu,
model dua fluida memperlakukan kedua fase sebagai kontinum yang saling menembus, digandeng oleh fraksi
volume dan pertukaran momentum.

\section{Di mana machine learning masuk}

Bab ini memperlihatkan tiga titik lemah simulasi aliran. Model penutup, seperti viskositas turbulen
dan model sub-grid, dibangun dari asumsi yang sering tidak berlaku. Simulasi yang cukup teliti mahal
sekali. Dan setiap kasus baru, dengan geometri atau syarat batas berbeda, harus disimulasikan dari awal.
Machine learning masuk lewat ketiga pintu itu: model penutup yang dipelajari dari data simulasi teliti
(\cref{ch:hibrida}), model pengganti yang menggantikan solver (\cref{ch:pinn,ch:operator}), dan metode
yang memanfaatkan pola berulang agar simulasi panjang bisa diekstrapolasi dengan murah (\cref{ch:rcfd}).

\sumberbab{Bab ini mengikuti kuliah Numerical Methods in Fluid Mechanics bagian volume hingga,
kopling tekanan dan kecepatan, turbulensi, panas dan spesies, serta multifase. Bacaan pendamping:
\citet{versteeg2007}, \citet{ferziger2020}, \citet{patankar1980}, dan \citet{pope2000}.}
